\subsection{The Dot Product}


\content{
\begin{definition} (The Dot Product) If $\vec{a}=\langle a_1,a_2,a_3\rangle$
and $\vec{b}=\langle b_1,b_2,b_3\rangle$, then the dot product
$\vec{a}\cdot\vec{b}$ is given by 
$$ \vec{a}\cdot \vec{b} = a_1b_1+a_2b_2+a_3b_3. $$
\end{definition}
\begin{example} if $\vec{a} =\langle 1, 0, 3\rangle$ and 
$\vec{b} =\langle 2,-1,1\rangle$, compute $\vec{a}\cdot \vec{b}$.
\end{example}
}{% presentation
\vspace{1in}
}{% nopresentation
\vspace{1in}
}{% comments
}


\content{
\begin{theorem} (Properties of the Dot Product) if $\vec{a}$, $\vec{b}$, and 
$\vec{c}$ are vectors in $\R^3$ and $c$ is a scalar, then 
\begin{enumerate}
\item $\vec{a}\cdot \vec{a} = |\vec{a}|^2$
\item $\vec{a}\cdot\vec{b} = \vec{b}\cdot \vec{a}$
\item $\vec{a}\cdot(\vec{b}+\vec{c}) = \vec{a}\cdot \vec{b} +
\vec{a}\cdot\vec{c}$
\item $(c\vec{a})\cdot \vec{b} =c(\vec{a}\cdot \vec{b})=\vec{a}\cdot
(c\vec{b})$
\item $\vec{0}\cdot \vec{a} = 0$
\end{enumerate}
\end{theorem}
}{% presentation
\vspace{2in}
}{% nopresentation
\vspace{1in}
}{% comments
Maybe provide a proof of 1 and 3.
}

\content{
\begin{theorem} If $\theta$ is the angle between vectors $\vec{a}$ and
$\vec{b}$, then 
$$ \vec{a}\cdot \vec{b} = |\vec{a}||\vec{b}|\cos\theta. $$
\end{theorem}
}{% presentation
\vspace{1in}
}{% nopresentation
\vspace{1in}
}{% comments
Prove this quickly. Draw a triangle with sides $\vec{a}$ and $\vec{b}$ and apply
the law of cosines: $c^2 = a^2 + b^2 -2ab\cos(\theta)$, where
$c=|\vec{a}-\vec{b}|$.
}

\content{
\begin{example} Find the angle between the vectors $\vec{a} = \la 2,2,-1 \ra$
and $\vec{b} = \la 5,-3,2\ra $.
\end{example}
}{% presentation
\vspace{2in}
}{% nopresentation
\vspace{1in}
}{% comments
}

\content{
\begin{theorem} Two vectors are orthogonal (perpendicular) if their dot product
is zero. 
\end{theorem}
}{% presentation
\vspace{2in}
}{% nopresentation
\vspace{2in}
}{% comments
Either prove this using the formula above, or use the Pythagorean Theorem. 
}

\subsubsection{Direction Angles}

\content{
\begin{definition}
The direction angles of a nonzero vector $\vec{a}$ are the angles $\alpha$,
$\beta$, $\gamma$ (all in the interval $[0,\pi]$) that $\vec{a}$ makes with the
positive $x$, $y$, and $z$-axes, respectively. 
\end{definition}
}{% presentation
\vspace{2in}
}{% nopresentation
\vspace{1in}
}{% comments
Draw the picture of an angle in the first octant.  Use the angle formula above
with the vectors $i$, $j$, and $k$ to derive formulas for the cosines of these
angles. Notice that squaring and summing these formulas gives an equality.
}

\content{
\begin{example}
Find the direction angles of the vector $\vec{a}=\la 1,2,3\ra$.
\end{example}
}{% presentation
\vspace{1in}
}{% nopresentation
\vspace{1in}
}{% comments
}

\subsubsection{Projections}

\content{
We will often need to break vectors up into components to determine how much
force or work is being done in a certain direction.  More specifically, if we
are given two vectors $\vec{a}$ and $\vec{b}$, we wish to determine the
projection of $\vec{a}$ onto $\vec{b}$. 
}{% presentation
\vspace{2in}
}{% nopresentation
\vspace{1in}
}{% comments
Draw both vectors $\vec{a}$ and $\vec{b}$ from the same starting point, then
drop the perpendicular from $\vec{a}$ onto $\vec{b}$ and derive the projection
formula from there. 

\textbf{Derive both the component formula and the projection formula. }
}

\content{
\begin{example} Find the projection of the vector $\vec{a}=\la 1,1,2\ra $ onto
the vector $\vec{b}=\la -2,3,1\ra$.
\end{example}
}{% presentation
\vspace{2in}
}{% nopresentation
\vspace{1in}
}{% comments
$$ \la -\frac{3}{7}, \frac{9}{14}, \frac{3}{14}\ra $$
}

\content{
\begin{example} One use of projections is to find the amount of work done by a
constant force. A wagon is being pulled a distance of 100 m along by a constant
force of 70 N. The handle of the wagon is held at an angle of $35^\circ$ above
the horizontal. Find the work done by the force. 
\end{example}
}{% presentation
}{% nopresentation
}{% comments
Set up a diagram, then when the triangle is drawn, use the fact that work is
force times distance to derive the formula 
$$ W = \vec{F}\cdot \vec{D}$$ 
where $\vec{F}$ is the force vector being applied and $\vec{D}$ is the
displacement vector.
}
